% unicode-math with STIX Two (bundle)
\documentclass{article}
\usepackage{amsmath}
\usepackage{fontspec}
\usepackage{unicode-math}
\setmainfont{STIX Two Text}
\setmathfont{STIX Two Math}
\begin{document}
\section{STIX Two}
Typography exists to honour content. Like oratory, music, dance, calligraphy -- like anything that lends its grace to language -- typography is an art that can be deliberately misused. It is a craft by which the meanings of a text (or its absence of meaning) can be clarified, honoured and shared, or knowingly disguised. In a world rife with unsolicited messages, typography must often draw attention to itself before it will be read.
\begin{equation}
  \hat{H}\psi = i\hbar\frac{\partial\psi}{\partial t}, \qquad
  \oint_{\partial\Sigma} \mathbf{B}\cdot d\boldsymbol{\ell} = \mu_0 I.
\end{equation}
\[ \left\lbrace \begin{aligned} x+y &= 3\\ x-y &= 1 \end{aligned}\right. \implies
   (x,y) = (2,1), \qquad \prod_{p\ \text{prime}} \frac{1}{1-p^{-s}} = \zeta(s). \]
\end{document}
